\documentclass[11pt,a4paper]{article}

\usepackage[margin=1.6cm]{geometry}
\usepackage{booktabs}
\usepackage{xcolor}
\usepackage{graphicx}
\usepackage{enumitem}
\usepackage[hidelinks]{hyperref}

\setlist{nosep,leftmargin=*}
\setlength{\parskip}{0.35em}
\setlength{\parindent}{0pt}

\title{\vspace{-1.5em}LiveRecall: Revision-Based Caption Distribution\\
\large CS3103 Assignment 4 Report\\
\large Team 23}
\author{\small Ngo Hoang Nhat Minh (A0301185N) \\ \small Xinyang (A0281296W) \and \small Chang Jing Jie Jordon (A0281173H) \\ \small Jihan (A0287033E)}
\date{}

\begin{document}
\maketitle
\vspace{-2em}

% =====================================================================
\section{Introduction and Architecture}

LiveRecall distributes a continuously revised caption stream to many
concurrent clients while preserving correctness, freshness and bounded
resource usage under slow readers, disconnections and constrained
networks. This report describes our server design, the revision and
recovery semantics it guarantees, the overload policy that keeps slow
clients from degrading normal ones, and the measurements that justify
both.

\textbf{Our Design priority:} Correctness is prioritized above freshness. Final Caption Completeness below 99\% caps the performance
assessment regardless of latency, and the specification treats
unbounded state as a correctness failure rather than a performance one.
Every bound in the system is therefore explicit and justified.

\subsection{Module structure}

\texttt{solution\_server.py} contains only wiring for CLI parsing,
logging configuration and connection dispatch. All behaviour are separated into
modules, which let four people work concurrently without breaking each other's work.

\begin{table}[h]
\centering\small
\begin{tabular}{@{}lll@{}}
\toprule
\textbf{Module} & \textbf{Responsibility} & \textbf{Key bound} \\
\midrule
\texttt{model.py}          & Record types, revision ordering rule & --- \\
\texttt{app.py}            & Configuration, dependency container  & all bounds and thresholds \\
\texttt{caption\_state.py} & Per-segment state, final durability  & 512 segments \\
\texttt{history.py}        & Recovery window, \textsc{history} replies & 90\,s / 1500 items \\
\texttt{hub.py}            & Client registry, fan-out             & --- \\
\texttt{client\_session.py}& Per-client queue, writer, overload policy & 32 soft / 128 hard frames; 3 in flight \\
\texttt{ingest.py}         & Parsing, validation, source handler  & --- \\
\texttt{metrics.py}        & Counters and observability           & --- \\
\bottomrule
\end{tabular}
\caption{Server modules. \texttt{model.py} and \texttt{app.py} are shared
interfaces; every other module has exactly one owner.}
\end{table}


\subsection{Data flow for one caption revision}

A caption revision follows this path:

\begin{enumerate}
\item \textbf{Validate} (\texttt{ingest}): incoming captions are ingested and invalidated. All untrusted input stops
      at this boundary.
\item \textbf{Admit} (\texttt{caption\_state.apply}): returns the
      update only if it advances state, and \texttt{None} for anything
      stale or duplicate. Rejected revisions therefore cannot reach the
      recovery window or any client queue.
\item \textbf{Record} (\texttt{history.record\_final}): a final caption enters
      the recovery window \emph{before} publication, so a client
      reconnecting within the same event-loop can read the caption history.
\item \textbf{Publish} (\texttt{hub.publish}): send the caption frame to
      every client queue synchronously, never awaiting a socket.
\end{enumerate}

Step~4 is why one unresponsive client cannot stall the broadcast path.
The naive baseline awaited every \texttt{send()} on the source's own
task; here the source only ever appends to bounded queues, and a
dedicated writer task per client drains to its own socket.

% =====================================================================
\section{Protocol and Revision Semantics}

We implement the mandatory compatibility interface without extensions,
which eliminates evaluator-compatibility risk. Messages are
\textsc{source\_hello}, \textsc{caption\_update},
\textsc{client\_hello} and \textsc{history}.

\subsection{Three invariants}

Revision handling must maintainthree rules (invariant) in this priority
order:

\begin{enumerate}
\item \textbf{Final durability.} Once a segment is final, nothing
      replaces it, even if it is a higher revision or a retransmission. This
      check precedes the revision comparison, so revision 99 still
      loses to a final revision 3.
\item \textbf{Monotonic revisions.} Within a segment a higher revision
      replaces a lower one; a lower one is never applied.
\item \textbf{Duplicate suppression.} A revision already held is
      rejected and not republished.
\end{enumerate}

Because all three are enforced at one admission point that returns
\texttt{None} on rejection, no downstream component re-implements
ordering.

% =====================================================================
\section{Recovery and Bounded State}

\subsection{Resume semantics}

\texttt{last\_segment}/\texttt{last\_revision} are interpreted as a
\textbf{watermark} which is the latest \emph{finalised} state the client already
holds. \textsc{history} returns only finals strictly newer than that
watermark, and no older than the requested window. If a client holds segment~12 but
missed segment~10, its watermark hides the gap and no bounded history
can repair it. The overload policy must never discard
a final, and the provisional revisions are the only droppable traffic.

\subsection{Why the window is bounded twice}

The recovery window is bounded by age (90\,s) \emph{and} by item count
(1500). The count bound is because a bursty source can
place far more than the expected number of segments inside a
90-second window, so a purely time-based bound would still grow with
the source rate. Live per-segment state is separately capped at 512
tracked segments, with an eviction watermark that keeps evicted
segments rejected so eviction cannot resurrect stale revisions, plus a
hard backstop for a source that abandons a segment without finalising
it.

\subsection{Measured recovery behaviour}

Against a running server, with 16 segments finalised:

\begin{table}[h]
\centering\small
\begin{tabular}{@{}llp{6.4cm}@{}}
\toprule
\textbf{Client state} & \textbf{Items returned} & \textbf{Interpretation} \\
\midrule
$(0,0)$, 30\,s  & 16, segments 1--16  & Late join receives the whole window \\
$(5,3)$, 30\,s  & 11, segments 6--16  & Watermark correctly excludes 1--5 \\
$(0,0)$, 5\,s   & 7,  segments 10--16 & Time cutoff is applied \\
\bottomrule
\end{tabular}
\caption{\textsc{history} contents for different resume states.}
\end{table}

% =====================================================================
\section{Backpressure and Overload Policy}

\subsection{A negative result that shaped the design}

We first triggered the overload policy on application queue depth, but with
10\,\% slow readers lagging by about 13\,s, peak queue depth was only
\textbf{1}. Probing a 500\,ms reader showed why: its first 38 frames were
acknowledged by pongs within 2\,ms while it was already falling behind, and
the first late pong (frame 40) took 5.8\,s. The client's WebSocket library
buffers 16 frames before pausing, and kernel and transport buffers absorb
more, so \emph{queue depth is not a usable lag signal}.

\subsection{The policy we adopted}

\textbf{Lag signal.} Each session pings its own connection. A ping shares the
caption byte stream, so its pong returns only after the client has read
everything sent before it; RTT above the connection's minimum is the queueing
delay in front of that client, and added network delay alone does not count.
Browsers and \texttt{websockets} answer pings automatically, so the protocol
is unchanged.

\textbf{Two modes.} In \emph{normal} mode every revision is sent in order. A
client becomes \emph{overloaded} when delay exceeds
\texttt{overload\_lag\_high} or its queue exceeds the soft bound, and clears
only after delay stays below \texttt{overload\_lag\_low}, with a short queue
and no coalescing, for \texttt{overload\_clear\_hold} seconds. While
overloaded:

\begin{itemize}
\item \textbf{Coalescing.} A newer revision overwrites a queued older
      provisional of the same segment in place; a final overwrites its own
      pending provisional. The queue holds at most one frame per active
      segment plus pending finals. Finals are never dropped; a client whose
      queue still exceeds the hard bound after coalescing is closed with code
      1013 and resumes through \textsc{history}.
\item \textbf{In-flight window.} At most \texttt{overload\_inflight\_frames}
      unconfirmed frames, so the backlog waits in our queue, where it can be
      coalesced, not in kernel buffers.
\item \textbf{Rate pacing.} The client's read rate is estimated from frames
      confirmed between two late pongs, when its queue is at the same level
      (within about 6\,\% of the true 2 messages/s). Sends are paced at
      0.5$\times$ that rate after a late pong, then 0.85$\times$, rising
      10\,\% every 5\,s without one, so a recovered client returns to full
      updates.
\end{itemize}

\begin{table}[h]
\centering\small
\begin{tabular}{@{}lll@{}}
\toprule
\textbf{Parameter} & \textbf{Value} & \textbf{Reason} \\
\midrule
\texttt{overload\_lag\_high}        & 0.25\,s & half the 500\,ms freshness threshold \\
\texttt{overload\_lag\_low}         & 0.08\,s & hysteresis gap \\
\texttt{overload\_clear\_hold}      & 2\,s    & four probe intervals \\
\texttt{overload\_inflight\_frames} & 3       & keeps backlog in our queue \\
\texttt{probe\_interval}            & 0.5\,s  & $\le$2 pings/s per healthy client \\
\bottomrule
\end{tabular}
\caption{Overload-policy parameters.}
\end{table}


\subsection{Why this is a correctness mechanism}

In T3 a 500\,ms reader consumes about 118 of the 228 revisions emitted in
60\,s. Spending that budget on finals and newest state rather than arrival
order raises Final Caption Completeness from 95.13\,\% to 99.34\,\%, across
the 99\,\% gate (\S\ref{sec:beforeafter}).




% =====================================================================
\section{Experiments}

All runs use the supplied harness with unmodified measurement logic, on
Ubuntu Linux (kernel 6.6.87.2, 4 vCPUs, 7.8\,GiB RAM, Python 3.13.5).
Each reported figure is the
median of three independent runs; \texttt{run\_demo.py --runs 3}
restarts the server for every run.

\subsection{Results}

\begin{table}[h]
\centering\small
\begin{tabular}{@{}lrrrrrr@{}}
\toprule
\textbf{Test} & \textbf{P50} & \textbf{P95} & \textbf{P99} &
\textbf{Fresh} & \textbf{Complete} & \textbf{Coverage} \\
 & (ms) & (ms) & (ms) & (\%) & (\%) & (\%) \\
\midrule
T1 20 clients, 20\,s          & 4.17 & 12.13 & 14.99 & 100.00 & 100.00 & 100.00 \\
T2 100 clients, 40\,s         & 11.38 & 29.39 & 36.66 & 100.00 & 100.00 & 100.00 \\
T3 100 clients, 10\,\% slow   & 14.48 & 2265.89 & 8414.06 & 94.65 & 99.34 & 95.18 \\
T4 100 clients, 20\,\% recon. & 10.58 & 25.78 & 30.95 & 100.00 & 100.00 & 98.33 \\
T5 delay+jitter+3\,\% loss    & 375.82 & 1928.09 & 6750.50 & 67.29 & 97.06 & 58.63 \\
\bottomrule
\end{tabular}
\caption{Public test results. Values are medians of three independent runs.
P50/P95/P99 are caption age in milliseconds; Fresh is the fraction at or
below 500\,ms. T5 applies 120\,ms mean delay, 30\,ms jitter and 3\,\% loss
on the loopback interface.}
\end{table}

Duplicate and stale caption counts were \textbf{zero} across every run,
as were invalid non-final \textsc{history} items.

\subsection{Slow-reader isolation (T3)}

\begin{table}[h]
\centering\small
\begin{tabular}{@{}lrrr@{}}
\toprule
\textbf{Population} & \textbf{P50 (ms)} & \textbf{P95 (ms)} & \textbf{Fresh (\%)} \\
\midrule
Normal clients (90\,\%) & 13.88    & 28.87   & 100.00 \\
Slow clients (10\,\%)   & 5550.85 & 9357.85 & 1.69 \\
\bottomrule
\end{tabular}
\caption{Normal clients are unaffected by deliberately slow readers.}
\end{table}

The server log shows the policy engaging exactly where intended:
\texttt{overload\_engaged}=10, matching the 10 slow readers, with no normal
client affected; \texttt{provisionals\_coalesced}=1020; and peak queue depth
9, against a soft bound of 32.

\subsection{Reconnect and recovery (T4)}

With 20\,\% of clients disconnected for 5\,s, final completeness
remained 100.00\,\% with zero duplicate or stale captions, and the
\textsc{history} response arrived with a median of 4.37\,ms and P95 of
7.97\,ms. This is measured
from the completed WebSocket handshake to \textsc{history} receipt, so it
covers the \textsc{client\_hello} round trip plus history construction;
the bounded window and indexed lookup keep this recovery path within a
single-digit millisecond tail in this experiment.

\subsection{Network impairment (T5)}

With \texttt{delay\_loss.sh lo} active, the three-run medians were
375.82\,ms P50, 1928.09\,ms P95 and 6750.50\,ms P99. Fresh Caption Ratio
fell to 67.29\,\%, while final completeness remained 97.06\,\%; duplicate,
stale and connection-error counts remained zero. WebSocket runs over TCP,
so packet loss does not normally lose application messages: retransmission
and head-of-line blocking instead inflate caption age, especially in the
tail. The completeness reduction reflects captions that remained
undelivered when the fixed experiment ended, rather than duplicate or
stale state.

\subsection{Before/after comparison}
\label{sec:beforeafter}

We compare the final design against the naive baseline (commit
\texttt{4c7e5c9}: a single shared broadcast via \texttt{asyncio.gather},
no recovery history, no per-client queues), using the same unchanged
harness and T3 parameters. Each figure is the median of three runs.

\begin{table}[h]
\centering\small
\begin{tabular}{@{}lrrl@{}}
\toprule
\textbf{Metric (T3 parameters)} & \textbf{Baseline} & \textbf{Final} & \textbf{Change} \\
\midrule
Final completeness (\%)              & 95.13   & 99.34   & +4.21 pp \\
Slow-client P50 / P95 age (s)        & 14.77 / 29.18 & 5.55 / 9.36 & $-$62.4\,\% / $-$67.9\,\% \\
Overall P99 age (ms)                 & 25120.57   & 8414.06    & $-$66.5\,\% \\
Normal-client P95 (ms)               & 18.78    & 28.87    & +10.09 ms \\
Slow-client freshness (\%)           & 1.69    & 1.69    & unchanged \\
Normal update coverage (\%)          & 95.18   & 95.18   & unchanged \\
\bottomrule
\end{tabular}
\caption{Before/after under identical parameters.}
\end{table}

\textbf{Trade-off.} The overload policy raises completeness and cuts slow-client
lag by roughly two thirds, but suppressing superseded provisionals reduces
coverage for the affected clients. Coverage returns to 100\,\% once no client is
lagging: T2 reports 100.00\,\% in all three runs, and in T3 the 90 normal clients
received all 20\,520 revisions. This is why the policy is triggered by measured
lag rather than applied unconditionally. The cost is a 10.09\,ms rise in
normal-client P95, most likely from per-client writer tasks and ping probes.


% =====================================================================
\section{Failures, Limitations and Future Work}
\label{sec:limits}

\subsection{Failures}

\begin{itemize}
\item \textbf{Queue depth was the wrong lag signal} (\S4.1): the backlog sits
      in client-library and kernel buffers, invisible to the server, so we
      switched to ping RTT above the minimum.
\item \textbf{Arrival-order delivery wastes slow readers' budget} on obsolete
      provisionals; coalescing fixed it (\S4.3).
\end{itemize}


\subsection{Limitations}

\begin{itemize}
\item A client offline for more than 90\,s has missed
      finals that no bounded history can return. This is the direct
      cost of refusing to retain an unbounded transcript.
\item Slow readers still lag by about 5\,s and miss the last few finals of a
      run, which are still in transit when the test ends. The client library
      buffers 16 frames before any server-side signal can see lag, so
      detection always costs about $16\times500$\,ms; removing it needs
      application-level acknowledgements, which the fixed protocol lacks.
\item One active session, as permitted by the
      specification. Multi-session routing is not implemented.
\item Currently, network impairment increases tail age. Under T5, P99 age
      rose to 6.75\,s and fresh-caption ratio fell to 67.29\,\%, while
      duplicate and stale counts stayed at zero. This is the expected
      reliable-TCP trade-off under loss and head-of-line blocking.
\end{itemize}


\subsection{Future work}
The main remaining cost is detection delay: the server only sees a slow
reader once its 16-frame library queue is full. An optional protocol
extension for client credits or acknowledgements would let the server
observe consumption directly instead of inferring it from late pongs,
removing that delay and improving slow-client freshness.

A second improvement would be adaptive history sizing based on observed
reconnect gaps, subject to a hard memory cap, rather than a fixed
90-second/1,500-item history.

% =====================================================================
\newpage
\section*{Appendix A: Reproduction}

\begin{verbatim}
python -m venv venv && ./venv/bin/pip install -r requirements.txt
python -m unittest discover -s tests -t .
ss -ltnp | grep 8765          # must be empty before a measured run
python run_demo.py --clients 100 --duration 60 --runs 3 \
    --slow-percent 10 --slow-ms 500 --output-stem t3
\end{verbatim}

Reported commands:
\begin{verbatim}
python -m unittest discover -s tests -t .
python run_demo.py --clients 20 --duration 20 --runs 3 --output-stem t1
python run_demo.py --clients 100 --duration 40 --runs 3 --output-stem t2
python run_demo.py --clients 100 --duration 60 --runs 3 \
    --slow-percent 10 --slow-ms 500 --output-stem t3
python run_demo.py --clients 100 --duration 60 --runs 3 \
    --disconnect-percent 20 --disconnect-at 15 --disconnect-seconds 5 \
    --output-stem t4
sudo ./network_profiles/delay_loss.sh lo
python run_demo.py --clients 100 --duration 40 --runs 3 --output-stem t5
sudo ./network_profiles/reset.sh lo
\end{verbatim}

\end{document}
